\documentclass[11pt]{article}

% useful packages
\usepackage[english]{babel}
\usepackage{amsmath}
\usepackage{amssymb}
\usepackage{wasysym}
\usepackage{mathtools}
\usepackage{floatpag}
\usepackage{cite}
\usepackage{hyperref}
\usepackage{verbatim}
\usepackage{listings}
\usepackage{xcolor}
\usepackage{graphicx}
\usepackage{caption}
\usepackage{mwe}
\usepackage{float}


% colorful listings

\lstset{language=C,%
    breaklines=true,
    keywordstyle=\color{blue},
    morekeywords=[2]{1}, keywordstyle=[2]{\color{black}},
    identifierstyle=\color{black},%
    stringstyle=\color{olive},
    commentstyle=\color{gray},%
    showstringspaces=false,%without this there will be a symbol in the places where there is a space
    numbers=left,
    numberstyle={\tiny \color{black}},% size of the numbers
    numbersep=9pt, % this defines how far the numbers are from the text
    emph=[1]{for,end,break,if},emphstyle=[1]\color{red}, %some words to emphasise
}




\begin{document}

\title{Solution to Homework 3}

\author{
   Meshach Aruoriwo Ogbor,\\
   Md Rasel}

\maketitle

\section{Exercise 3.1 - Point to Point Communication}
\lstinputlisting[language=C]{ex3_1.c}
%For compilation we use gcc 11.4 and openMPI 4.1.2. We compile with\\
%\verb+ mpicc -o hello_mpi hello_mpi.c+
For compilation on Stromboli, we used the following command: \\
\verb+ mpicc -o ex3_1 ex3_1.c+
\newpage

\section{Submit Script for exercise 3.1}
\lstinputlisting[language=bash]{job.sh}
\section{Results}
\verbatiminput{output.txt}
%Our main findings are summarized by fig.~\ref{fig1}.

%\begin{figure}
%   \includegraphics[width=\linewidth]{plot.pdf}
%   \caption{Skill level over time}\label{fig1}
%\end{figure}

\newpage
\section{Exercise 3.2 - Tree}
\lstinputlisting[language=C]{ex3_2.c}
%For compilation we use gcc 11.4 and openMPI 4.1.2. We compile with\\
%\verb+ mpicc -o hello_mpi hello_mpi.c+
For compilation on Stromboli, we used the following command: \\
\verb+ mpicc -o ex3_2 ex3_2.c -lm+ 
\newpage

\section{Submit Script for exercise 3.2}
\lstinputlisting[language=bash]{job3_2.sh}
\section{Results}
\verbatiminput{output3_2.txt}
As you can see from the results, time taken by the custom ``Tree'' function is less than the custom ``Loop'' function. For example: on rank 0, it took around 4 times less time (0.123695/0.032437) while using the ``Tree'' function. Theoretically, we should get $\log_2 48 = 5.5$ times less time. 
\end{document}
